\section{Secant--hyperplane moments}
\label{sec:moments}

The character equations~\eqref{eq:character} count hyperplanes through
zero, one, and two cap points.  Restricting them to the hyperplanes
through a fixed secant leaves the same three counts, but the pairs of
further cap points now split according to whether they are coplanar with
the secant, and that is where \(T_\ell\) enters.

\begin{proof}[Proof of Theorem~\ref{thm:secant-hyperplane-moments}]
Fix the secant \(\ell=ab\).  The hyperplanes through \(\ell\) correspond
to the hyperplanes of the quotient space \(\PG(d-2,q)\), so there are
\(\theta_{d-2}\) of them.  A further cap point \(u\) spans with \(\ell\) a
plane, which lies in \(\theta_{d-3}\) hyperplanes; summing over the
\(k-2\) choices of \(u\) gives the first moment.  For a pair
\(\{u,v\}\subseteq A\setminus\{a,b\}\), the span of \(\ell,u,v\) is a
plane when \(\{a,b,u,v\}\) is coplanar and a solid otherwise, since no
three of the four points are collinear.  A plane lies in
\(\theta_{d-3}\) hyperplanes and a solid in \(\theta_{d-4}\); there are
\(T_\ell\) coplanar pairs and \(\binom{k-2}{2}-T_\ell\) others, so
\[
 \sum_{\Pi\supset\ell}\binomtwo{w_\Pi}
 =\Bigl(\binomtwo{k-2}-T_\ell\Bigr)\theta_{d-4}+T_\ell\,\theta_{d-3}
 =\binomtwo{k-2}\theta_{d-4}+q^{d-3}\,T_\ell,
\]
because \(\theta_{d-3}-\theta_{d-4}=q^{d-3}\).
\end{proof}

For \(d=3\) the hyperplanes through \(\ell\) are its planes,
\(\theta_{-1}=0\), and the second moment is \(T_\ell\) itself, which is
the pencil identity~\eqref{eq:pencil}.  For \(d\ge4\) the theorem
expresses \(T_\ell\) through the section sizes of the hyperplanes through
\(\ell\),
\begin{equation}
 T_\ell=\frac1{q^{d-3}}\Bigl(\sum_{\Pi\supset\ell}\binomtwo{w_\Pi}-\binomtwo{k-2}\theta_{d-4}\Bigr),
\label{eq:T-from-sections}
\end{equation}
so the sum on the right is congruent to \(\binom{k-2}{2}\theta_{d-4}\)
modulo \(q^{d-3}\) for every secant.

\begin{remark}[Shortened codes]
In the dictionary of Section~\ref{sec:preliminaries}, the hyperplanes
through \(\ell=ab\) are the codewords of \(C_A\) vanishing at the
coordinates \(a\) and \(b\), that is, the code \(C_A\) shortened at
\(\{a,b\}\), a \([k-2,d-1]_q\) code whose codewords have weights
\(k-2-w_\Pi\).  Its dual is \(C_A^\perp\) punctured at \(\{a,b\}\), whose
words of weight at most two are the \((q-1)T_\ell\) images of the
weight-four words of \(C_A^\perp\) with support containing \(a\) and
\(b\).  Theorem~\ref{thm:secant-hyperplane-moments} is therefore the
first two power moment identities of Pless~\cite{Pless1963} for the
shortened code, which express its zeroth, first, and second weight
moments through its length, dimension, and the numbers of dual words of
weight one and two.  We keep the geometric proof because it exhibits the
plane-versus-solid dichotomy that produces the factor \(q^{d-3}\).
\end{remark}

The same restriction can be made to the hyperplanes through any fixed
subspace; counting over the hyperplanes through a fixed plane is the
technique of Thas~\cite{Thas2017} for bounding the largest caps of
\(\PG(n,q)\), \(q\) even.  The secant is the right subspace here because
the second moment then measures the collisions of one secant with the
others, the quantity that the coverage bound of
Section~\ref{sec:coverage} consumes.
